\chapter{Hướng dẫn sử dụng}

\section{Khởi động chương trình}
\begin{enumerate}
    \item Đảm bảo đã biên dịch dự án bằng hệ thống \texttt{CMake}. Thư mục \texttt{Release} sẽ chứa các file thực thi (xem Phụ lục A).
    \item Chạy \texttt{build\_features.exe} để chương trình duyệt và nén toàn bộ kho ảnh thực tế vào các file cơ sở dữ liệu \texttt{.yml}.
    \item Chạy \texttt{cbir\_gui.exe} để khởi động màn hình làm việc chính của hệ thống. 
\end{enumerate}

\section{Thực hiện truy vấn}
\begin{enumerate}
    \item Tại cửa sổ dòng lệnh, chọn ảnh bằng cách nhập đường dẫn hoặc ứng dụng sẽ lấy ngẫu nhiên 1 file trong thư mục Test.
    \item Trên giao diện đồ họa, bấm vào các nút thuật toán ở thanh công cụ bên dưới (Ví dụ: ColorHist, SIFT, HOG...).
    \item Bấm vào nút \textbf{K=3}, \textbf{K=5}, \textbf{K=11}, \textbf{K=21} để tùy chọn số lượng hiển thị ảnh mong muốn.
    \item Nhấn chuột trái để xác nhận truy vấn. Hệ thống sẽ tự động đối sánh và xuất danh sách ảnh.
    \item Quan sát hình ảnh được trả về trực tiếp trên cửa sổ ứng dụng và xem thêm thời gian tính toán thực tế.
\end{enumerate}

\missingfigure{main_ui.png}{Giao diện chính của chương trình.}

\section{Ý nghĩa thông tin trên giao diện}
\begin{table}[H]
\centering
\begin{tabularx}{\textwidth}{|l|X|}
\hline
\textbf{Thành phần} & \textbf{Ý nghĩa} \\
\hline
Ảnh lớn bên trái & Ảnh truy vấn gốc (Query Image). \\
\hline
Lưới ảnh bên phải & Bảng danh sách các ảnh gần giống nhất do hệ thống nhận diện. \\
\hline
Chữ \texttt{Dist} / \texttt{Sim} & Điểm số tương đồng hoặc khoảng cách (càng nhỏ càng khớp nếu là Dist). \\
\hline
Chữ \texttt{Time} & Thời gian tiêu tốn cho phiên truy vấn vừa xong (đơn vị: ms). \\
\hline
Danh sách nút bấm & Quản lý các Extractor đang được nạp và số lượng ảnh. Nút đổi màu khi kích hoạt. \\
\hline
\end{tabularx}
\caption{Chú giải chi tiết giao diện CBIR.}
\end{table}
